\section{Resultados}
Una vez con esto procedimos a usar las metricas SD1 y SD2 para hacer una clasificacion de
nuestros datos, primero, obtuvimos los SD1 y SD2 para cada sujeto y cada nodo dado que son una
gran cantidad de datos a continuacion solo se presentan la informacion de los primeros nodos.

\begin{table}[H]
\centering
\caption{Valores de SD1 por sujeto control - Nodos C3 a F4}
\label{tab:sd1_nodos}
\begin{tabular}{lSSSSS}
\toprule
\textbf{Persona} & \textbf{C3} & \textbf{C4} & \textbf{CZ} & \textbf{F3} & \textbf{F4} \\
\midrule
EMV & 0.8826 & 0.6504 & 0.5307 & 2.0511 & 1.4118 \\
GH2 & 1.1255 & 0.7941 & 0.6204 & 1.2813 & 0.8536 \\
GUR & 0.8111 & 0.7287 & 0.5889 & 1.2971 & 1.0394 \\
JAL & 0.1724 & 0.1760 & 0.1525 & 0.1917 & 0.1722 \\
JAN & 0.2843 & 0.5670 & 0.2766 & 0.3238 & 0.8579 \\
MGN & 0.1631 & 0.1550 & 0.1559 & 0.2360 & 0.3016 \\
MJN & 0.2275 & 0.3874 & 0.1925 & 0.3752 & 0.3050 \\
MMA & 0.3375 & 0.3271 & 0.2239 & 0.3374 & 0.3861 \\
RAN & 0.6604 & 0.4555 & 0.3940 & 0.5738 & 0.6582 \\
VCN & 3.2163 & 0.6690 & 0.9445 & 3.6138 & 1.2387 \\
\bottomrule
\end{tabular}
\end{table}

\begin{table}[H]
\centering
\caption{Valores de SD2 por sujeto - Nodos C3 a F4}
\label{tab:sd2_nodos}
\begin{tabular}{lSSSSS}
\toprule
\textbf{Persona} & \textbf{C3} & \textbf{C4} & \textbf{CZ} & \textbf{F3} & \textbf{F4} \\
\midrule
EMV & 2.6623 & 2.7021 & 2.7963 & 3.9621 & 5.2372 \\
GH2 & 2.9317 & 2.8607 & 2.7682 & 4.0477 & 4.2576 \\
GUR & 8.2427 & 8.2315 & 8.2305 & 8.6304 & 8.7627 \\
JAL & 2.2067 & 2.4451 & 2.7602 & 4.2273 & 4.8202 \\
JAN & 3.0633 & 3.2070 & 3.1277 & 4.0145 & 3.8344 \\
MGN & 3.1882 & 3.5294 & 2.9507 & 3.7230 & 6.7881 \\
MJN & 3.0583 & 2.7824 & 2.4490 & 4.1814 & 3.7262 \\
MMA & 1.7083 & 1.6141 & 1.8008 & 1.9580 & 2.0338 \\
RAN & 2.4456 & 3.2833 & 2.7073 & 3.3289 & 5.7529 \\
VCN & 2.7899 & 1.7255 & 1.7045 & 3.4208 & 2.1108 \\
\bottomrule
\end{tabular}
\end{table}

\begin{table}[H]
\centering
\caption{Valores de SD1 por sujeto DCL - Nodos C3 a F4}
\label{tab:sd1_dcl_nodos}
\begin{tabular}{lSSSSS}
\toprule
\textbf{Persona} & \textbf{C3} & \textbf{C4} & \textbf{CZ} & \textbf{F3} & \textbf{F4} \\
\midrule
AEF & 0.2349 & 0.2883 & 0.2153 & 0.3263 & 0.4077 \\
CLM & 0.2887 & 0.3060 & 0.2414 & 0.2835 & 0.2897 \\
FGV & 0.3798 & 0.3797 & 0.3471 & 0.5096 & 0.7053 \\
JGM & 0.7672 & 0.9304 & 0.4838 & 1.0032 & 1.2345 \\
LIV & 0.3173 & 0.3777 & 0.3363 & 0.4066 & 0.3927 \\
PCM & 1.0591 & 0.9398 & 0.6573 & 1.6471 & 1.5755 \\
RLM & 0.2198 & 0.2246 & 0.1929 & 0.3291 & 0.3886 \\
RRM & 0.7946 & 0.8183 & 0.6009 & 0.7816 & 0.7177 \\
\bottomrule
\end{tabular}
\end{table}

\begin{table}[H]
\centering
\caption{Valores de SD2 por sujeto DCL - Nodos C3 a F4}
\label{tab:sd2_dcl_nodos}
\begin{tabular}{lSSSSS}
\toprule
\textbf{Persona} & \textbf{C3} & \textbf{C4} & \textbf{CZ} & \textbf{F3} & \textbf{F4} \\
\midrule
AEF & 2.0202 & 1.6376 & 1.8530 & 3.2175 & 3.1491 \\
CLM & 2.4997 & 2.9559 & 2.7775 & 3.6257 & 4.3958 \\
FGV & 2.3190 & 1.8103 & 2.1587 & 3.3382 & 3.0057 \\
JGM & 5.4872 & 4.7901 & 4.3497 & 5.9940 & 5.8812 \\
LIV & 5.3111 & 4.8494 & 4.9797 & 5.9089 & 6.0851 \\
PCM & 4.8881 & 4.9062 & 3.9286 & 6.9373 & 7.2490 \\
RLM & 5.2707 & 5.0991 & 5.2509 & 6.2569 & 5.8319 \\
RRM & 2.0499 & 2.3600 & 2.0627 & 2.9606 & 2.8267 \\
\bottomrule
\end{tabular}
\end{table}

Ahora usando el KNN hay dos formas de proceder una vez obtuvimos las metricas SD1 y SD2
la primera es clasificando usando la media de SD1 y SD2 por cada sujeto y usando un solo sujeto
como prueba y el resto como entrenamiento, o bien usar todos los valores de SD1 y SD2, es decir
considerar cada valor de SD1 y SD2 de cada nodo y de cada persona, para el primer caso tenemos
lo siguiente.
Consederando k = 3 el modelo KNN obtuvo una presiciòn total de 0,667 es decir clasifico 12
sujetos bien del total de 18 sujetos de la siguiente manera:

\begin{table}[H]
\centering
\caption{Clasificación de sujetos - Real vs Predicho}
\label{tab:clasificacion}
\begin{tabular}{lcc}
\toprule
\textbf{Sujeto} & \textbf{Real} & \textbf{Predicho} \\
\midrule
EMV  & Control & DCL \\
GH2  & Control & DCL \\
GUR  & Control & DCL \\
JAL  & Control & Control \\
JAN  & Control & Control \\
MGN  & Control & Control \\
MJN  & Control & Control \\
MMA  & Control & DCL \\
RAN  & Control & Control \\
VCN  & Control & DCL \\
AEF  & DCL     & DCL \\
CLM  & DCL     & Control \\
FGV  & DCL     & DCL \\
JGM  & DCL     & DCL \\
LIV  & DCL     & DCL \\
PCM  & DCL     & DCL \\
RLM  & DCL     & DCL \\
RRM  & DCL     & DCL \\
\bottomrule
\end{tabular}
\end{table}

De donde obtenemos la matriz de confusiòn:

\begin{table}[H]
\centering
\caption{Clasificación de sujetos - Real vs Predicho (Random Forest)}
\label{tab:clasificacion1_rf}
\begin{tabular}{lcc}
\toprule
\textbf{} & \textbf{Control} & \textbf{DCl} \\
\midrule
Control & 5 & 5 \\
DCL  & 1 & 7 \\
\bottomrule
\end{tabular}
\end{table}

Ahora si procedemos con la segunda opcion obtenemos la siguiente informacion:

\begin{table}[H]
\centering
\caption{Clasificación de sujetos - Real vs Predicho}
\label{tab:clasificacion2}
\begin{tabular}{lcc}
\toprule
\textbf{Sujeto} & \textbf{Real} & \textbf{Predicho} \\
\midrule
EMV  & Control & Control \\
GH2  & Control & Control \\
GUR  & Control & DCL \\
JAL  & Control & DCL \\
JAN  & Control & Control \\
MGN  & Control & DCL \\
MJN  & Control & Control \\
MMA  & Control & DCL \\
RAN  & Control & Control \\
VCN  & Control & Control \\
AEF  & DCL     & Control \\
CLM  & DCL     & Control \\
FGV  & DCL     & Control \\
JGM  & DCL     & Control \\
LIV  & DCL     & DCL \\
PCM  & DCL     & Control \\
RLM  & DCL     & DCL \\
RRM  & DCL     & Control \\
\bottomrule
\end{tabular}
\end{table}

De donde obtenemos la matriz de confusiòn:

\begin{table}[H]
\centering
\caption{Clasificación de sujetos - Real vs Predicho (Random Forest)}
\label{tab:clasificacion2_rf}
\begin{tabular}{lcc}
\toprule
\textbf{} & \textbf{Control} & \textbf{DCl} \\
\midrule
Control & 6 & 4 \\
DCL  & 7 & 1 \\
\bottomrule
\end{tabular}
\end{table}


En esta opcion obtuvimos que la presiciòn total del modelo es de solo 0,444 es decir de los 18
sujetos a clasificar solamente clasificò bien 8, lo cual nos indica que bajo las metricas que trabajamos
es mejor considerar las medias de SD1 y SD2 para obtener una mayor presiciòn.
Ahora con el algoritmo de random forest los resultados que se obtuvieron no fueron relevantes
sin embargo sirve para hacer una buena comparativa y analizar la situaciòn.
De igual forma considerando solamente las medias para SD1 y SD2 se obtuvo lo siguiente:

\begin{table}[H]
\centering
\caption{Clasificación de sujetos - Real vs Predicho (Random Forest)}
\label{tab:clasificacion_rf}
\begin{tabular}{lcc}
\toprule
\textbf{Sujeto} & \textbf{Real} & \textbf{Predicho} \\
\midrule
EMV  & Control & Control \\
GH2  & Control & DCL \\
GUR  & Control & DCL \\
JAL  & Control & Control \\
JAN  & Control & Control \\
MGN  & Control & Control \\
MJN  & Control & Control \\
MMA  & Control & DCL \\
RAN  & Control & Control \\
VCN  & Control & DCL \\
AEF  & DCL     & Control \\
CLM  & DCL     & Control \\
FGV  & DCL     & Control \\
JGM  & DCL     & Control \\
LIV  & DCL     & DCL \\
PCM  & DCL     & Control \\
RLM  & DCL     & Control \\
RRM  & DCL     & Control \\
\bottomrule
\end{tabular}
\end{table}

De donde obtuvimos la matriz de confusiòn siguiente

\begin{table}[H]
\centering
\caption{Clasificación de sujetos - Real vs Predicho (Random Forest)}
\label{tab:clasificacion5_rf}
\begin{tabular}{lcc}
\toprule
\textbf{} & \textbf{Control} & \textbf{DCl} \\
\midrule
Control & 6 & 4 \\
DCL  & 7 & 1 \\
\bottomrule
\end{tabular}
\end{table}

De donde se obtuvo una presiciòn total del 0,389 es decir de los 18 sujetos solamente clasificò
bien a 7.
Con el otro caso, es decir, usando todos los nodos se obtuvo:

\begin{table}[H]
\centering
\caption{Clasificación de sujetos - Real vs Predicho (Random Forest)}
\label{tab:clasificacion_rf2}
\begin{tabular}{lcc}
\toprule
\textbf{Sujeto} & \textbf{Real} & \textbf{Predicho} \\
\midrule
EMV  & Control & Control \\
GH2  & Control & DCL \\
GUR  & Control & DCL \\
JAL  & Control & Control \\
JAN  & Control & DCL \\
MGN  & Control & Control \\
MJN  & Control & Control \\
MMA  & Control & DCL \\
RAN  & Control & DCL \\
VCN  & Control & Control \\
AEF  & DCL     & Control \\
CLM  & DCL     & Control \\
FGV  & DCL     & Control \\
JGM  & DCL     & DCL \\
LIV  & DCL     & DCL \\
PCM  & DCL     & Control \\
RLM  & DCL     & Control \\
RRM  & DCL     & Control \\
\bottomrule
\end{tabular}
\end{table}

De donde obtuvimos la matriz de confusiòn siguiente

\begin{table}[H]
\centering
\caption{Clasificación de sujetos - Real vs Predicho (Random Forest)}
\label{tab:clasificacion4_rf}
\begin{tabular}{lcc}
\toprule
\textbf{} & \textbf{Control} & \textbf{DCl} \\
\midrule
Control & 5 & 5 \\
DCL  & 6 & 2 \\
\bottomrule
\end{tabular}
\end{table}

Dando una presiciòn total del 0.389, es decir de igual forma solo clasificò bien 7 sujetos de los
18 que hay.




